% Part: computability
% Chapter: computability-theory
% Section: reducibility

\documentclass[../../../include/open-logic-section]{subfiles}

\begin{document}

\olfileid[fr]{cmp}{thy}{red}
\olsection{Réductibilité}

\begin{explain}
Nous savons maintenant qu'au moins un ensemble,
$K_0$, est effectivement énumérable sans être
calculable. Il en existe d'autres. La méthode des
réductions permet de montrer que d'autres ensembles
ont ces propriétés, sans reprendre chaque fois
les définitions de départ.

De façon générale, une « réduction » de l'ensemble
$A$ à l'ensemble~$B$ transforme une réponse à la
question de savoir si des !!{element}s appartiennent
à~$B$ en une réponse à la question correspondante
pour~$A$. Nous étudierons la \emph{réduction many-one},
mais il existe d'autres notions de réduction, aux
propriétés diverses. Elles jouent aussi un rôle
central en complexité des calculs, où l'efficacité
entre en jeu. Par exemple, un ensemble est dit
« NP-complet » s'il appartient à NP et si tout
problème de NP s'y réduit par une notion semblable
à celle définie ci-dessous, mais qui exige en plus
que la réduction soit calculable en temps polynomial.

Nous avons déjà employé implicitement une réduction.
Définissons
\[
K = \Setabs{x}{\cfind{x}(x) \fdefined},
\]
soit $K=\Setabs{x}{x\in W_x}$. La preuve de
l'indécidabilité du problème de l'arrêt
(\olref[hlt]{thm:halting-problem}) montre directement
que $K$ n'est pas calculable. Rappelons que
\[
K_0 = \Setabs{\tuple{e, x}}{\cfind{e}(x) \fdefined },
\]
soit $K_0=\Setabs{\tuple{e,x}}{x\in W_e}$.\footnote{Correction éditoriale : la reformulation source inverse la paire en $\tuple{x,e}$. Nous conservons l'ordre $\tuple{e,x}$ de la définition affichée : indice du programme, puis entrée.}
De toute preuve de la non-calculabilité de~$K$,
on déduit facilement celle de~$K_0$ : si $K_0$
était calculable, on pourrait décider si
!!a{element}~$x$ appartient à~$K$ en demandant
si la paire $\tuple{x,x}$ appartient à~$K_0$.
La fonction~$f$ définie par
$f(x)=\tuple{x,x}$ est une \emph{réduction}
de $K$ à~$K_0$.
\end{explain}

\begin{defn}
Soient $A$ et~$B$ des ensembles d'entiers naturels.
Une fonction calculable $f\colon\Nat\to\Nat$ est une
\emph{réduction many-one} de $A$ à~$B$ si, pour tout
entier naturel~$x$,
\[
x \in A \quad \text{si et seulement si} \quad f(x) \in B.
\]
S'il existe une telle fonction~$f$, on dit que $A$
est \emph{many-one réductible} à~$B$, et l'on écrit
$A\leq_m B$. Si chaque ensemble est many-one
réductible à l'autre, on dit que $A$ et~$B$ sont
\emph{many-one équivalents}, et l'on écrit
$A\equiv_m B$.
\end{defn}

Si la fonction~$f$ de cette définition est en outre
injective, on parle de \emph{réduction injective}
(ou « one-one ») de $A$ à~$B$. La plupart des
réductions données ci-dessous possèdent cette
propriété plus forte, que nous n'utiliserons pas.

\begin{digress}
Il est vrai, mais nullement évident, que la
réductibilité injective est strictement plus
forte que la réductibilité many-one. Autrement
dit, il existe des ensembles infinis $A$ et~$B$
tels que $A$ soit many-one réductible à~$B$,
mais non réductible à~$B$ par une fonction injective.
\end{digress}

\end{document}
